\begin{textsample}{POS Dim 1 – pi – Score 60.00 – pi/pi\_em\_40.txt}  \label{ex:f1_pos_054}
Área de Matemática e suas Tecnologias Parte Comum do Currículo Alinhado à BNCC e a DCN do Ensino Médio Nos últimos anos do \textbf{século} XIX, a educação passou por grandes transformações.
Estas transformações se deram por grandes movimentos, como, por exemplo, a Escola Nova e a Escola Ativa, que vieram dar às instituições de ensino um viés \textbf{transformador}, mudando o seu modo de ser através de \textbf{métodos} ativos, onde integrar o \textbf{aluno} ao ato de “ fazer ” se torna elemento essencial no processo de ensino e aprendizagem.
Muito antes disso, na Antiguidade, as pessoas se reuniam em várias situações para conversar, discutir, trocar ideias, e sem perceberem, umas ensinavam aquilo que sabiam às outras, de forma prática, experimental e investigativa.
E embora tais \textbf{métodos} e movimentos tenham surgido para traçar um novo painel ao ensino propedêutico, o que se vê na maioria das escolas é o distanciamento cada vez maior entre o que se deve ensinar e o que se deve aprender, tornando a educação escolar enfadonha e desgastada, segundo a visão dos estudantes.
A Matemática, uma das mais questionadas pela metodologia pouco atrativa, vem sendo ensinada sem se relacionar com a realidade e expectativa dos \textbf{alunos}, não \textbf{instigando} nenhuma curiosidade naquilo que o professor \textbf{diz} estar ensinando, isto é, o ensino da \textbf{disciplina} afasta os estudantes, ao invés de \textbf{instigar} a paixão pelos números.
Em \textbf{suma}, não estamos preparando nossos jovens e adultos para a revolução \textbf{digital}, pois o futuro será daqueles que saibam Matemática suficiente para se integrar nesse processo.
Infelizmente, a manipulação de poder ou a \textbf{abordagem} didática, muitas vezes exercida pelo professor durante a \textbf{aula}, conduz claramente à não percepção do que os silêncios na sala de \textbf{aula} indicam, a ausência do diálogo no contexto escolar, e isso é um dos motivos do fracasso do ensino de Matemática.
Segundo Baraldi, “ para os \textbf{alunos}, a Matemática consiste num manipular de fórmulas que, após \textbf{certo} treino, torna -se fácil em situações próprias da Matemática ” ( Baraldi, 1999, p. 88 ).
No entanto, a Matemática não é um cardápio de fórmulas em que se tenha que decorá -las para resolver um \textbf{problema}, mas, sobretudo, uma forma de pensar e enxergar o mundo que nos \textbf{auxilia} no dia a dia.
Educar em Matemática requer objetivos estabelecidos em função de objetos do conhecimento, planejamento para a ação educativa e ferramentas para potencializar o ensino e, não menos importante, a avaliação daquilo que foi realizado.
É preciso eliminar a concepção dos que pensam a sala de \textbf{aula} como lugar de encontro de \textbf{alunos} ignorantes, que têm ali um professor totalmente sábio repassando informações puramente técnicas, sem vida própria e distante da realidade.
É notório que compreendamos a sala de \textbf{aula} como um espaço de múltiplas interações de pessoas com conhecimentos de senso comum e que ali estão almejando a aquisição \textbf{sistematizada} de conhecimentos.
Para Chamay ( 2001, p. 36 ), “ um dos objetivos essenciais do ensino da Matemática é precisamente que, o que se ensine, esteja carregado de significado, tenha sentido para o \textbf{aluno} ”.
O sentido de um conhecimento matemático se define não só pela coleção de situações em que este conhecimento é realizado como \textbf{teoria} matemática ou como situações em que o sujeito o encontrou como meio de solução, mas também pelo conjunto de concepções que rejeita, de \textbf{erros} que evita, de economias que procura, de formulações que retorna ( Bousseau, 1983 ).
O objeto do conhecimento matemático, quando compreendido, tem por significado inicial ter a consciência de que ele tem, em algum lugar, um tipo de existência, e que tal existência é algo dependente de aspectos históricos e sociais.
Assim, a Matemática deve ser inicialmente observada como ela se apresenta, levando à \textbf{compreensão} de que a investigação e a descoberta, ao longo dos anos, se constituíram como essenciais a ela, bem como seus núcleos inspiradores : um mundo externo com seus \textbf{problemas} e sua própria estruturação interna.
Sobre isso, Cortella ( 1998, p. 102 ) afirma que : > Quando um educador(a) nega ( com ou sem intenção ) aos \textbf{alunos} a \textbf{compreensão} das condições culturais, históricas e sociais de produção do conhecimento, termina por reforçar a mitificação e a sensação de perplexidade, impotência e incapacidade cognitiva.
O ensino proposto pela escola do \textbf{século} XXI desafia formar pessoas solidárias, criativas e autônomas, capazes de lidar em meio às incertezas, buscando alcançar uma sociedade bem mais justa, digna e solidária.
O ensino de Matemática é bem maior do que apenas provocar o \textbf{raciocínio} lógico-dedutivo dos \textbf{alunos}, que por muitas vezes parece ser o objetivo principal em planejamentos de \textbf{aulas}, pois através do conhecimento da história da Matemática é possível a \textbf{compreensão} do presente, o entendimento do passado e a \textbf{projeção} de futuro.
O Lúdico também é parte importante nesse processo, conectando o ensino com experiências adequadas a cada faixa etária.
Além disso, o uso de \textbf{tecnologias} se incorpora ao processo cognitivo como parte indissociável e \textbf{indispensável} para a formação na sociedade da Comunicação e da \textbf{Tecnologia} da atualidade.
Segundo os \textbf{Parâmetros} Curriculares Nacionais ( 1997 ) : > ( … ) A Matemática é componente importante na construção da cidadania, na medida em que a sociedade utiliza, cada vez mais, de conhecimentos científicos e recursos tecnológicos, dos quais os cidadãos devem se apropriar.
A aprendizagem em Matemática está ligada à \textbf{compreensão}, isto é, à apreensão do significado ; aprender o significado de um objeto ou acontecimento pressupõe vê -lo em suas relações com outros objetos e acontecimentos.
Recursos didáticos como jogos, livros, \textbf{vídeos}, calculadora, computadores e outros materiais têm um papel importante no processo de ensino aprendizagem.
Contudo, eles precisam estar integrados a situações que levem ao exercício da \textbf{análise} e da reflexão, em última \textbf{instância}, a base da atividade matemática.
Com a BNCC, a Matemática para o Ensino Médio é apresentada com \textbf{foco} na construção de uma visão integrada, aplicada à realidade, nos diversos contextos, levando em consideração as vivências cotidianas dos \textbf{discentes}, estes que são inteiramente \textbf{impactados} pelas revoluções tecnológicas, pelo que \textbf{exige} o \textbf{mercado} de trabalho, pelos projetos de vida dos povos, entre outros.
Assim, a área da Matemática e suas Tecnologias no Ensino Médio tem a importante responsabilidade de aproveitar os conhecimentos adquiridos pelos estudantes durante o Ensino Fundamental, possibilitando o avanço no \textbf{letramento} matemático aprendido nos anos anteriores de escolaridade.
Segundo a nova proposta, aprender Matemática deve consistir, essencialmente, em fazer Matemática.
E fazer Matemática na escola significa tratar com questões que \textbf{exigem} observação, \textbf{análise}, estabelecimento de conexões, conjecturas, percepção e expressão de regularidades, busca de explicações, criação de soluções, invenção de estratégias próprias que envolvam noções, conceitos e \textbf{métodos} matemáticos e, por fim, a comunicação da produção realizada.
Argumentação e \textbf{prova} são dois aspectos da comunicação em Matemática.
Na sociedade \textbf{moderna} e \textbf{complexa} que \textbf{vivemos} hoje, ensinar Matemática requer o desenvolvimento de habilidades básicas, através de uma aprendizagem significativa e que evidencie sua \textbf{aplicabilidade}, que dê as possibilidades reais de conquistas.
A Educação Matemática deve visar a construção de um saber para capacitar nossos \textbf{alunos} a pensar e a refletir sobre a realidade, assim como a agir e transformá -la.
Santaló ( 2001, pg. 11 ) afirma ainda que : > A missão dos educadores é, portanto, preparar as novas gerações para o mundo em que terão que viver.
Isto quer \textbf{dizer} proporcionar -lhes o ensino necessário para que adquiram as destrezas e habilidades que vão necessitar para seu desempenho, com comodidade e eficiência, no seio da sociedade que enfrentarão ao concluir sua escolaridade.
Nesse aspecto, o ensino de Matemática prestará sua contribuição à medida que forem exploradas metodologias que priorizem a criação de estratégias, a comprovação, a justificativa, a argumentação, o espírito \textbf{crítico}, e favoreçam a criatividade, o trabalho coletivo, a iniciativa pessoal e a autonomia advinda do desenvolvimento da confiança na própria \textbf{capacidade} de conhecer e enfrentar desafios ( PCN’s – Matemática, 1997 ).
Enquanto isso, a melhoria da qualidade da educação escolar no Brasil se firma, na atualidade, como um dos grandes desafios a serem superados pelas gestões públicas e \textbf{pesquisadores} da educação, dada a condição pouco satisfatória da qualidade da aprendizagem apresentada pela maioria dos estudantes ao \textbf{final} da Educação Básica, sejam estes oriundos de escolas públicas ou privadas.
A \textbf{prova} do PISA – Programa Internacional de Avaliação de Estudantes, coordenada pela OCDE – Organização para a Cooperação e o Desenvolvimento Econômico, da qual o Brasil participa como convidado desde 2000, demonstra a real situação do país em \textbf{termos} de aprendizagem, já que tem o objetivo de avaliar de forma amostral se os estudantes de escolas públicas e particulares, aos 15 anos de idade, adquiriram conhecimentos e habilidades essenciais para uma participação plena em sociedade \textbf{moderna}.
Ao avaliar a \textbf{proficiência} dos estudantes de 15 anos de idade em Leitura, Matemática e Ciências de escolas públicas e privadas, que estejam matriculados a partir do 8º ano do Ensino Fundamental, o objetivo principal do PISA é fornecer indicadores que contribuam para a discussão da qualidade da educação ofertada nos países participantes, de modo a \textbf{subsidiar} políticas de melhoria da educação básica.
Por outro \textbf{lado}, os resultados produzidos pelo programa podem indicar até que ponto as escolas de cada país participante estão preparando seus jovens para exercer e/ou responder às \textbf{demandas} sociais, já que o PISA avalia até que ponto os estudantes, próximos ao \textbf{final} da educação obrigatória, adquiriram conhecimentos e habilidades essenciais para plena participação na vida social e \textbf{econômica}.
Conforme Pessoa e Silva ( 2019 ), o resultado brasileiro nesta avaliação comparada, isto é, promovida entre os países e convidados da OCDE, indica que os estudantes brasileiros avaliados em 2012, cujo \textbf{foco} da avaliação era Matemática, possuem baixos níveis de aprendizado matemático, considerando que cerca de 70 \% encontram -se nos níveis menores da escala, conforme a seguir : | Nível/Edição | Abaixo de 1 | 1 | 2 | 3 | 4 | 5 | 6 | |---|---:|---:|---:|---:|---:|---:|---:| | 2012 | 35 \% | 32 \% | 21 \% | 8 \% | 4 \% | 1 \% | 0,5 \% | Tabela 1 – Resultado Brasil 2012 – Matemática Fonte : INEP ( 2013 ).
Para as \textbf{autoras}, esta constatação é validada ao observarmos a diferença entre o Brasil e países que obtêm os melhores resultados no PISA, \textbf{verificada} nas etapas posteriores da vida escolar.
Segundo a OCDE, a média de engenheiros formados entre os concluintes de cursos superiores nos países membros da organização é de 14 \% ; na Coreia do Sul, esse índice é de 25 \%.
No Brasil, segundo dados do MEC/INEP, obtidos através do Censo da Educação Superior ( 2016 ), apenas cerca de 6 \% dos concluintes estavam nas áreas de Engenharia.
A distância observada na condução das políticas educacionais é ainda mais esmagadora se considerarmos o desenvolvimento tecnológico e \textbf{econômico} observado entre os dois países.
Nesse ínterim, é, portanto, válido considerar através desses percentuais a vocação e o incentivo que os dois países fornecem para a inovação tecnológica.
Segundo o boletim Radar : \textbf{tecnologia}, produção e comércio exterior do Instituto de Pesquisa Econômica e Aplicada – IPEA ( 2014 ) : > [ … ] a economia nacional mostra -se pouco intensiva em trabalho qualificado e de cunho técnico-científico, insumos tidos como necessários à inovação e à competição em \textbf{mercados} globais.
De outro, o sistema educacional brasileiro parece \textbf{operar} em um equilíbrio de baixa qualidade, em todos os níveis de ensino, o que afeta a \textbf{capacidade} produtiva da força de trabalho, em geral, e a das engenharias, em particular ( RADAR, 2014, p. 19 ).
A formação em nível superior nas carreiras que \textbf{exigem} \textbf{aprofundamento} científico fica comprometida pela pouca eficiência dos sistemas de ensino ofertados nacionalmente e, nesse âmbito, “ o Brasil \textbf{segue} tendo diante de si grandes desafios para melhorar o aprendizado, em todos os níveis de escolaridade, com \textbf{potenciais} repercussões na qualidade da força de trabalho ” ( RADAR, p. 34 ).
Em âmbito estadual, o Piauí não foge à regra brasileira : desde 2006 recebendo nota padronizada no teste, o Estado obteve resultados mais baixos que a média nacional, embora com pequenos avanços.
Conforme dados do PISA/INEP ( 2016 ), o Piauí concentra a maioria dos \textbf{alunos} ( 79 \% ) nos níveis abaixo de 1 e 1 em Matemática.
Tais números revelam que o padrão de desempenho dos estudantes piauienses demonstra carência de aprendizagem em relação ao que é previsto para a etapa de escolaridade avaliada, tendo em \textbf{vista} que a OCDE institui o nível 2 como o “ aceitável ” para os estudantes com 15 anos de idade.
Dessa forma, o resultado do Piauí na avaliação corrobora a \textbf{fragilidade} do sistema educacional brasileiro diante das desigualdades \textbf{econômicas} e sociais historicamente \textbf{verificadas} e debatidas em âmbito nacional, se incorporando às diferentes concepções educacionais observadas nas regiões brasileiras.
Tradicionalmente, os estados da região Norte e Nordeste detêm, em média, os \textbf{piores} indicadores educacionais e sociais ; no PISA, essa realidade não é diferente.
Tais desigualdades educacionais estão muito presentes no cenário brasileiro.
Na tentativa de amenizar essas irregularidades, a Base Nacional Comum Curricular – BNCC se consolida como uma oportunidade para a transformação dos sujeitos, a partir da aquisição de novas \textbf{competências}.
A reformulação curricular proposta com a BNCC tem o objetivo de \textbf{aproximar} mais a escola das expectativas de aprendizagem dos estudantes do \textbf{século} XXI, que são altamente tecnológicos e que buscam respostas para suas \textbf{demandas} \textbf{socioemocionais}, evidenciando o protagonismo juvenil/estudantil, bem como sua preparação para o mundo do trabalho.

% matched lemmas: abordagem, aluno, análise, aplicabilidade, aprofundamento, aproximar, aula, autor, auxiliar, capacidade, certo, competência, complexo, compreensão, crítico, demanda, digital, discente, disciplina, dizer, econômico, erro, exigir, final, foco, fragilidade, impactar, indispensável, instigar, instância, lado, letramento, mau, mercado, moderno, método, operar, parâmetro, pesquisador, potencial, problema, proficiência, projeção, prova, raciocínio, segar, sistematizar, socioemocional, subsidiar, suma, século, tecnologia, teoria, termos, transformador, verificar, vista, vivar, vídeo
\end{textsample}
